\chapter{Your own first person camera}

This is the chapter people brace themselves for, and it is smaller than it looks.
There are two angles and one formula, and then you are done.

\section{The two angles}

\begin{center}
\begin{tikzpicture}[scale=1.15,every node/.style={font=\sffamily\footnotesize}]
  \fill[rlight,rounded corners=3pt] (-2.4,-2.3) rectangle (3.2,2.3);

  % ground circle (top view of yaw)
  \draw[codeframe,fill=white] (0,0) circle (1.7);
  \draw[->,rred,line width=1pt] (0,0) -- (1.7,0) node[right,text=rred] {$+X$};
  \draw[->,rblue,line width=1pt] (0,0) -- (0,-1.7) node[below,text=rblue] {$+Z$};

  \draw[->,rink,line width=1.3pt] (0,0) -- (30:1.7) node[above right,text=rink] {forward};
  \draw[rorange,line width=1pt] (0.75,0) arc (0:30:0.75);
  \node[text=rorange] at (0.95,0.2) {yaw};

  \node[align=left,text=rgrey] at (0,2.0) {seen from above: \textbf{yaw} turns you around Y};
  \node[align=left,text=rgrey] at (0,-2.05) {\textbf{pitch} tilts you up and down (not shown)};
\end{tikzpicture}
\end{center}

\begin{itemize}[leftmargin=*,itemsep=2pt]
  \item \textbf{yaw} --- the horizontal turn, spinning around the Y axis. Looking
        left and right.
  \item \textbf{pitch} --- the vertical tilt, around the X axis. Looking up and
        down.
\end{itemize}

There is a third, \emph{roll}, which tilts your head sideways. First person games
almost never use it, so we will not either.

Both are stored in \textbf{radians}, because that is what C's \lstinline|sinf|
and \lstinline|cosf| want. A full turn is $2\pi$. Convert for display with
raylib's \lstinline|RAD2DEG| and \lstinline|DEG2RAD| constants.

\section{From angles to a direction}

\begin{lstlisting}
Vector3 forward = {
    cosf(pitch)*cosf(yaw),
    sinf(pitch),
    cosf(pitch)*sinf(yaw)
};
\end{lstlisting}

If trigonometry is a distant memory, here is the intuition and you can leave it
at that. Ignore pitch for a second: on the flat ground, \lstinline|cos(yaw)| and
\lstinline|sin(yaw)| trace a circle of radius 1 as yaw goes round. That is the
direction you are facing, on the floor.

Now add pitch. \lstinline|sin(pitch)| is how much you are looking up, and it goes
straight into Y. \lstinline|cos(pitch)| shrinks the flat part: look straight up
and nothing is left of the horizontal direction, which is exactly right.

Starting value: \lstinline|yaw = -PI/2| points you along $-Z$, which is
conventionally ``into the screen''. \lstinline|pitch = 0| is the horizon.

\section{From a direction to sideways}

To strafe you need the direction to your right. That is the \textbf{cross
product} of where you are looking and where up is:

\begin{lstlisting}
Vector3 right = Vector3Normalize(
    Vector3CrossProduct(forward, (Vector3){ 0.0f, 1.0f, 0.0f }));
\end{lstlisting}

All you need to know about the cross product: give it two directions and it hands
back a third one perpendicular to both. Forward and up are perpendicular-ish, and
the thing perpendicular to both of them is sideways. That is the entire reason it
appears in every camera in existence.

\section{Moving}

\begin{lstlisting}
Vector3 flatForward = Vector3Normalize((Vector3){ forward.x, 0.0f, forward.z });

Vector3 move = { 0 };
if (IsKeyDown(KEY_W)) move = Vector3Add(move, flatForward);
if (IsKeyDown(KEY_S)) move = Vector3Subtract(move, flatForward);
if (IsKeyDown(KEY_D)) move = Vector3Add(move, right);
if (IsKeyDown(KEY_A)) move = Vector3Subtract(move, right);

if (Vector3Length(move) > 0.0f) move = Vector3Normalize(move);

camera.position = Vector3Add(camera.position, Vector3Scale(move, speed*dt));
camera.target   = Vector3Add(camera.position, forward);
\end{lstlisting}

That last line is the one that ties it all together: the camera looks at a point
one unit in front of the eye, in the direction we computed. Forget it and the
camera moves but never turns.

\section{The four details that make or break the feel}

\begin{enumerate}[leftmargin=*,itemsep=5pt]
  \item \textbf{Flatten the forward vector for walking.} Notice
        \lstinline|flatForward| above: the Y component is thrown away and the
        result renormalized. Without it, looking at the sky and pressing W makes
        you fly.

  \item \textbf{Clamp the pitch} to just under 90 degrees:
\begin{lstlisting}
if (pitch >  1.55f) pitch =  1.55f;     // about 89 degrees
if (pitch < -1.55f) pitch = -1.55f;
\end{lstlisting}
        At exactly 90, \lstinline|forward| is parallel to \lstinline|up|, the
        cross product collapses to zero, \lstinline|right| becomes undefined and
        the camera snaps violently. This is \emph{gimbal lock}, it is famous, and
        one \lstinline|if| avoids it.

  \item \textbf{Normalize the movement vector.} Hold W and D together without it
        and the two unit vectors add up to a length of $\sqrt{2}$: you move 41\%
        faster diagonally. Players find that exploit within a minute.

  \item \textbf{Multiply by \lstinline|dt|.} As always.
\end{enumerate}

\section{Reading the mouse}

\begin{lstlisting}
DisableCursor();                       // once, at startup

Vector2 mouse = GetMouseDelta();
yaw   += mouse.x*0.003f;
pitch -= mouse.y*0.003f;               // note the MINUS
\end{lstlisting}

Pitch subtracts because screen Y grows downwards while looking up is positive
pitch --- the 2D and 3D Y axes disagreeing one final time. Flip that sign and you
have inverted mouse look, which some people actually prefer, so make it an
option rather than a bug.

The \lstinline|0.003f| is mouse sensitivity, in radians per pixel. Expose it as a
setting; everybody wants a different one.

\section{The finishing touch: head bob}

Pure decoration, three lines, and it is most of what makes movement feel
physical:

\begin{lstlisting}
if (walking) bobTimer += dt*speed*1.6f;
float bob = walking ? sinf(bobTimer)*0.045f : 0.0f;
camera.position.y = EYE_HEIGHT + bob;
\end{lstlisting}

A sine wave on the eye height, running faster when you run. Keep the amplitude
tiny --- 4 centimetres is plenty. Anything bigger and players feel seasick.

\tryit{Lesson 4 of the 3D course is this chapter as a running program, printing
yaw, pitch and the forward vector live so you can watch the numbers as you turn.}{./course3d/build.sh 4}
